\section{AI Agent 带来的新型攻击向量}

\subsection{Prompt Injection}

在 AI 系统中隐藏或注入误导性指令，使其偏离预期行为。这是 AI Agent 面临的最核心安全威胁。

\subsection{Tool Misuse}

通过欺骗性 prompt 操纵 Agent 滥用其集成的工具。例如，诱导 Agent 执行未授权的文件操作或 API 调用。

\subsection{Intent Breaking}

操纵 Agent 的执行计划（plan），将其行动从原始意图重定向到攻击者的目标。

\subsection{Identity Spoofing}

利用被攻破的认证机制冒充合法 Agent。

\subsection{Code Attacks}

利用 Agent 执行代码的能力，获取对执行环境的未授权访问。

\begin{warningbox}{AI Agent 的攻击面远大于传统软件}
传统软件的攻击面主要在输入验证和认证环节。AI Agent 的攻击面则扩展到了 prompt（可被注入）、工具调用（可被滥用）、执行计划（可被篡改）、身份认证（可被伪造）和代码执行环境（可被利用）等多个维度。
\end{warningbox}

\subsection{本章小结}

AI Agent 引入了至少五种新型攻击向量。防护需要在 prompt 过滤、工具权限控制、计划验证、身份验证和沙箱隔离等多个层面同时展开。

